\documentclass[conference]{IEEEtran}
\IEEEoverridecommandlockouts
\usepackage{cite}
\usepackage{amsmath,amssymb,amsfonts}
\usepackage{algorithmic}
\usepackage{graphicx}
\usepackage{textcomp}
\usepackage{xcolor}
\usepackage{booktabs}
\def\BibTeX{{\rm B\kern-.05em{\sc i\kern-.025em b}\kern-.08em
    T\kern-.1667em\lower.7ex\hbox{E}\kern-.125emX}}

% ---------------------------------------------------------------
% PLACEHOLDER MACRO
% Every \PH{...} marks information that is NOT yet verified from the
% experimental data or from the source PDFs. Resolve all of them
% before submission; search the file for "\PH{" to list them.
% ---------------------------------------------------------------
\newcommand{\PH}[1]{\textcolor{red}{[\,#1\,]}}

\begin{document}

\title{The Effect of XAI Explanation Types on User Trust\\ and Purchase Intention in a Conversational\\ Recommender System for the Perfume Domain}

\author{\IEEEauthorblockN{1\textsuperscript{st} Raissa Desyandita}
\IEEEauthorblockA{\textit{\PH{department}} \\
\textit{\PH{institution}}\\
\PH{City, Country} \\
\PH{email / ORCID}}
\and
\IEEEauthorblockN{2\textsuperscript{nd} \PH{Supervisor Name}}
\IEEEauthorblockA{\textit{\PH{department}} \\
\textit{\PH{institution}}\\
\PH{City, Country} \\
\PH{email / ORCID}}
}

\maketitle

\begin{abstract}
Explanations are widely assumed to increase user trust in AI-based recommender
systems, yet accumulating evidence indicates that the effect depends on the kind
of explanation provided rather than on its mere presence. Comparative evidence on
explanation types remains scarce in conversational recommender systems, and is
almost absent in sensory-hedonic domains, where the defining product attribute
cannot be transmitted through a screen. This study reports a between-subjects
experiment with 90 participants, 30 per condition, using a conversational
recommender system for perfume in which a large language model handles dialogue
while a knowledge-graph embedding backend produces the recommendations.
Explanation type was manipulated purely through system-prompt conditioning, with
no fine-tuning, across three conditions: feature-based, narrative-based, and
comparative-based. Perceived trust and purchase intention were measured with
adapted two-item scales and analysed using Kruskal-Wallis and Dunn tests.
Explanation type significantly affected both trust and purchase intention, with a
large and a medium-to-large effect size respectively. Contrary to the common
design intuition that warmer and more relatable explanations build trust,
narrative-based explanations produced the lowest trust, significantly below both
objective-grounded conditions. Comparative-based explanations produced the highest
purchase intention, significantly above the other two conditions. The association
between trust and purchase intention was positive but weak (Spearman rho = 0.229,
p = 0.030), and no evidence of moderation by perfume familiarity was found. The
two design goals were therefore not achieved by the same condition, which
indicates a concrete design trade-off for conversational recommender systems in
hedonic domains.
\end{abstract}

\begin{IEEEkeywords}
explainable AI, conversational recommender systems, perceived trust, purchase
intention, hedonic domain, user experience
\end{IEEEkeywords}

% ===============================================================
\section{Introduction}
% ===============================================================

Conversational recommender systems (CRS) depart from the one-shot interaction
paradigm of classical recommenders by supporting a task-oriented, multi-turn
dialogue in which preferences are elicited incrementally and suggestions can be
justified, questioned, and revised \cite{jannach2021}. The recent maturation of
large language models (LLMs) has made such systems inexpensive to build and
natural to use, and has correspondingly increased the practical importance of a
question that predates them: what should the system say when it explains itself?

Explainable AI (XAI) has been proposed as the bridge between opaque recommendation
models and users who must decide how far to rely on them \cite{rong2024}. The
canonical framing holds that explanations may serve seven distinct aims, including
transparency, effectiveness, persuasiveness, and trust, and, crucially, that these
aims can be mutually incompatible, so that an explanation optimised for one may
degrade another \cite{tintarev2012}. Empirical work in marketing supports the
first half of this framing: across four studies, the presence of post hoc
explanations raised interpretability perceptions and, through them, trust,
purchase intention, and click-through on AI recommendations \cite{chen2024}.

The second half of the framing, however, is frequently overlooked in practice.
Explanation presence alone does not guarantee calibrated trust. In an experiment
with 215 participants, feature-importance explanations produced higher trust than
counterfactual ones, and explanations raised trust even when the underlying system
was unreliable, which the authors interpret as overreliance rather than warranted
trust \cite{brito2023}. Related work shows that explanation interfaces can induce
trust calibration errors, in which users agree or disagree with the system without
deliberation \cite{naiseh2021}. Taken together, these results shift the research
question from whether an explanation is present to which kind of explanation is
given.

Two gaps follow from this shift, and this paper addresses both.

\textit{First, explanation types are rarely compared head-to-head inside a single
conversational system.} The seminal demonstration that different explanation types
build different trusting beliefs, with \textit{how} explanations raising competence
and benevolence, \textit{why} explanations raising benevolence, and
\textit{trade-off} explanations raising integrity, was conducted on a static
recommendation agent for digital cameras nearly two decades ago \cite{wang2007}.
More recent comparative work contrasts feature-importance with counterfactual
explanations \cite{brito2023}, or catalogues contrastive and counterfactual
generation methods \cite{stepin2021}, but does so outside a conversational
e-commerce setting. A comparison of feature-based, narrative-based, and
comparative-based justifications within one LLM-driven CRS has not, to our
knowledge, been reported.

\textit{Second, sensory-hedonic domains remain underexplored relative to
utilitarian ones.} This matters because the effect of explanation is not domain
invariant. Chen et al. found the facilitating effect of post hoc explanations to be
stronger in utilitarian than in hedonic decision domains, and further found that
explanation type interacts with domain: attribute-based explanations were more
effective in utilitarian settings, whereas user-based explanations were more
effective in hedonic ones \cite{chen2024}. Converging evidence from clickstream
data shows that negative reviews depress purchase probability for utilitarian
products but not for hedonic ones, and not when the review concerns taste-related
factors \cite{varga2024}. Product type likewise moderates the path from perceived
deception to online repurchase intention \cite{wangx2022}. Findings imported from
rational domains therefore cannot be assumed to transfer.

Perfume is an extreme instance of the hedonic case. Its defining attribute, scent,
cannot be transmitted digitally at all, unlike the visual attributes of clothing or
the auditory attributes of music. Its descriptive vocabulary, comprising terms such
as accord, sillage, and olfactory family, is mastered mainly by enthusiasts, which
makes domain familiarity a plausible moderator. Purchase decisions are strongly
identity-linked, and online purchase uncertainty is correspondingly high. If
explanation design has domain-dependent effects, perfume is where those effects
should be most visible.

Accordingly, this study asks three research questions.
\textbf{RQ1:} Does explanation type affect perceived trust and purchase intention?
\textbf{RQ2:} Is perceived trust associated with purchase intention?
\textbf{RQ3:} Does perfume familiarity moderate the effect of explanation type?

The contributions of this paper are as follows.

\begin{itemize}
\item \textbf{An artefact.} A perfume CRS in which an LLM handles language and a
knowledge-graph embedding backend produces recommendations, and in which the
explanation mode is switched by system-prompt conditioning alone. Because the
recommendation engine is untouched by the manipulation, the recommended items are
held constant across conditions, which isolates \textit{how the system explains}
from \textit{what the system recommends}.

\item \textbf{Comparative empirical evidence in a hedonic domain.} A
between-subjects experiment with 90 participants establishes that explanation type
significantly affects both trust and purchase intention in a domain whose core
attribute is not digitally transmissible, responding directly to the call for
domain-sensitive XAI evidence raised in \cite{chen2024}.

\item \textbf{A counter-intuitive finding with theoretical grounding.}
Narrative-based explanations, designed to be the warmest and most relatable,
produced the \textit{lowest} trust of the three conditions. This inverts the
prevailing design heuristic that humanising an agent increases confidence in it,
and is consistent with evidence that competent conversational styles outperform
warm ones on cognition-oriented tasks \cite{libaoku2024}.

\item \textbf{An explicit design trade-off.} The condition maximising trust and the
condition maximising purchase intention were not the same, and the trust-purchase
intention association was weak. This provides empirical support for the claim that
explanatory aims can conflict \cite{tintarev2012}, and suggests that comparative
explanations act on purchase intention through a route that does not pass through
trust.

\item \textbf{Transparently reported null and weak results.} No evidence of
moderation by familiarity was found, and the limitations of statistical power that
qualify this result are reported rather than suppressed.
\end{itemize}

The remainder of this paper is organised as follows. Section~\ref{sec:related}
reviews explanation types in XAI and recommender systems, explanations in CRS,
trust and purchase intention, and hedonic product domains, and positions the
present work. Section~\ref{sec:system} describes the system architecture and the
three explanation conditions. Section~\ref{sec:method} presents the experimental
design, measures, manipulation check, and analysis plan. Section~\ref{sec:results}
reports the results. Section~\ref{sec:discussion} interprets them and states the
limitations. Section~\ref{sec:conclusion} concludes.

% ===============================================================
\section{Related Work}
\label{sec:related}
% ===============================================================

\subsection{Explanation Types in XAI and Recommender Systems}

Tintarev and Masthoff established the reference taxonomy of explanatory aims for
recommender systems, comprising effectiveness, satisfaction, transparency,
scrutability, trust, persuasiveness, and efficiency, and argued that any evaluation
must state which aim it targets because the aims can be incompatible
\cite{tintarev2012}. Their own empirical studies illustrate the point sharply:
across two domains, personalising feature-based explanations was
\textit{detrimental} to effectiveness while it plausibly improved satisfaction, a
result the authors explicitly report as contrary to their hypothesis. The
structural lesson, that a more appealing explanation can be a less useful one, is
directly relevant to the present study.

Wang and Benbasat provided the foundational demonstration that explanation type,
not merely explanation presence, determines which component of trust is affected
\cite{wang2007}. Using a purpose-built recommendation agent, they showed that
\textit{how} explanations increased competence and benevolence beliefs,
\textit{why} explanations increased benevolence beliefs, and \textit{trade-off}
explanations increased integrity beliefs. Their conclusion, that the effects of
explanations depend largely on the content and type of explanation provided, is the
premise on which the present study rests.

Subsequent work has refined the picture without resolving it. Stepin et al. survey
contrastive and counterfactual explanation generation and propose a two-level
taxonomy, noting that contrastive explanations are in high demand yet lack
standardised evaluation \cite{stepin2021}. De Brito Duarte et al. compared feature
importance against counterfactual explanations in an AI recommendation task and
found feature importance to yield higher trust among non-expert participants, while
also observing that explanations increased trust in an unreliable system
\cite{brito2023}. They note that counterfactual explanations may nonetheless be
advantageous where expert knowledge is available, a caveat that becomes relevant to
the interpretation of our comparative condition. Naiseh et al. approach the same
tension from a design perspective, deriving five principles intended to support
trust calibration rather than trust maximisation \cite{naiseh2021}.

At the level of the field, Rong et al. reviewed 97 user studies of model
explanations and concluded that the effectiveness of XAI in users' interaction with
models is \textit{not consistent across applications}, and that user evaluations
remain sparse and rarely informed by cognitive or social science \cite{rong2024}.
Domain-specific studies confirm the value of measuring perception directly rather
than assuming it; de Campos et al., for instance, evaluate an explainable
content-based recommender against transparency, satisfaction, trust, and
scrutability in a user study \cite{decampos2024}.

The anchor study for the present work is Chen et al., who conducted four studies on
consumer responses to explainable AI recommendations \cite{chen2024}. Three findings
are pertinent. First, the presence of post hoc explanations increased
interpretability perceptions, which in turn increased trust, purchase intention, and
click-through. Second, this facilitating effect was stronger in the utilitarian than
in the hedonic decision-making domain. Third, explanation type moderated
effectiveness, with attribute-based explanations more effective in utilitarian
domains and user-based explanations more effective in hedonic domains. Their own
conclusion is that there are no universally optimal post hoc explanations. The
present study extends this line from static recommendation displays into multi-turn
conversation, and, as reported in Section~\ref{sec:results}, obtains a result that
runs against the direction their third finding would predict.

\subsection{Explanations in Conversational Recommender Systems}

Jannach et al. provide the standard survey of CRS and observe that although
explanations are often considered a main feature of a convincing dialogue, these
aspects are not explored a great deal \cite{jannach2021}. The intervening years have
partially, but only partially, closed this gap.

Hernandez-Bocanegra and Ziegler investigated conversational explanation delivery
directly, contrasting a GUI-based dialogue with a chatbot-like natural language
interface and varying the degree of interactivity \cite{ziegler2023}. Using
structural equation modelling, they related perceived explanation quality to
transparency, trust, and effectiveness, and found that offering interactive options
for scrutinising explanatory arguments significantly improved user evaluations. Their
work establishes that the conversational channel itself shapes how explanations are
received, but it manipulates interface and interactivity rather than explanation
content.

Siro et al. annotated CRS dialogues along six aspects, namely relevance,
interestingness, understanding, task completion, interest arousal, and efficiency,
and found that the concept of satisfaction varies across users, with the system's
ability to make relevant recommendations the dominant turn-level factor
\cite{siro2023}. Rana et al. studied critiquing-based conversational recommendation
with 228 participants and showed that the choice of underlying recommendation
algorithm materially changes critiquing behaviour, dialogue length, and satisfaction
\cite{rana2024}.

Methodologically, the closest parallel to the present study is Ziegfeld et al., who
built three versions of a CRS and ran a between-subjects experiment with 66
participants, measuring ten user experience constructs \cite{ziegfeld2025}. Notably,
they found \textit{no} significant effects of their manipulation on any
questionnaire-based subjective construct, although objective measures such as chat
duration and recommendation match score did differ. Their manipulation concerned
preference elicitation method rather than explanation content. The contrast with the
present findings is instructive and is revisited in Section~\ref{sec:discussion}.

\subsection{Trust and Purchase Intention in Recommender Systems}

Trust in a technological artefact is conventionally decomposed into beliefs about
competence, benevolence, and integrity \cite{mcknight2002}, and this
three-dimensional structure underlies the explanation-to-trust mapping reported in
\cite{wang2007}. For user-centric evaluation of recommenders more broadly, the ResQue
framework of Pu, Chen, and Hu supplies validated constructs and design guidance,
including the guidance that explaining \textit{why} an item is recommended increases
transparency whereas explaining \textit{trade-offs} increases persuasion
\cite{pu2012}. Perceived usefulness is measured following the technology acceptance
tradition \cite{davis1989}.

The presumed chain from trust to purchase behaviour has been validated in the chatbot
context by Yen and Chiang, who combined 204 questionnaire responses with EEG evidence
from 30 participants and found credibility, competence, anthropomorphism, social
presence, and informativeness to predict trust, which in turn predicted purchase
intention \cite{yen2021}. Notably, they also found that media richness and
playfulness did \textit{not} increase trust, reasoning that communication with a
shopping chatbot is a functional task in which decorative cues are not what users
seek. This is an early signal of the pattern the present study observes.

Work on anthropomorphism qualifies the chain further. Cheng et al. found that both
perceived warmth and perceived competence raise trust in chatbots, but that an
exchange relationship, that is, a transactional context, strengthens the weight of
competence \cite{cheng2022}. Li et al. modelled dual pathways in which chatbot warmth
and competence generate both trust and chatbot overload, with both routes feeding
service evaluation and subsequent purchase, which implies that trust is not the sole
channel to behavioural intention \cite{liyang2025}. Crolic et al. supplied the
strongest cautionary precedent, showing across five studies that anthropomorphising a
service chatbot \textit{lowers} satisfaction, firm evaluation, and purchase intention
for customers who arrive angry, via expectancy violation \cite{crolic2022}. The
proposition that a more human-seeming agent is a better agent is therefore already
contested in top-tier marketing research.

Most directly relevant, Li et al. examined anthropomorphic interaction style through
the lens of mind perception across three experiments \cite{libaoku2024}. They found
that for cognition-oriented service tasks a competent interaction style produced more
positive consumer responses than a warm one, mediated by agent-mind perception,
whereas for emotion-oriented tasks the reverse held, mediated by experience-mind
perception. Their central claim, that no style is universally superior and that task
type sets the boundary, provides the theoretical mechanism this paper invokes to
explain its counter-intuitive result.

\subsection{Sensory and Hedonic Product Domains}

Hedonic products are acquired for enjoyment, experience, and self-expression rather
than function, and are judged against criteria that are subjective, difficult to
verbalise, and unstable relative to the binary and stable criteria applied to
utilitarian goods \cite{wangx2022}. This difference is not merely descriptive; it
changes how information is processed.

Varga and Albuquerque, using clickstream data from a major online retailer, found
that a single negative review significantly reduces purchase probability and
increases search for substitutes, but that these effects hold for utilitarian product
categories and not for hedonic ones, and hold for reviews about functionality or
customer service but not for reviews about taste-related factors such as design,
colour, or material \cite{varga2024}. Wang et al. showed across three studies that
perceived deception damages online repurchase intention less for hedonic than for
utilitarian products, attributing the difference to the fluctuant and susceptible
nature of hedonic evaluation \cite{wangx2022}. Chen et al., as noted, found the effect
of post hoc explanations themselves to be domain-contingent \cite{chen2024}.

The cumulative implication is that XAI evidence gathered in rational domains does not
automatically transfer to sensory ones. Existing sensory work concentrates on food and
wine; to our knowledge no comparative test of XAI explanation types has been reported
in the olfactory domain.

\subsection{Familiarity as a Boundary Condition}

Domain familiarity is a plausible moderator of explanation effectiveness. Bayer et al.,
using a purpose-built decision support system in the chess domain with 100
participants, found that domain-specific expertise \textit{negatively} affects trust in
an AI-based system, and, more pointedly for the present design, that the effect of
explanation on trusting intention is moderated by expertise: users with high domain
knowledge trust the system more when an explanation is provided, whereas users with low
domain knowledge do not \cite{bayer2023}. Their conclusion is that explanations must fit
the user.

Celar and Byrne provide the closest design analogue, crossing explanation type with
domain familiarity explicitly across four experiments with 731 participants
\cite{celar2023}. They found counterfactual explanations to be judged more helpful than
causal ones, and, importantly for the comparative condition studied here, found that
counterfactual explanations improved the accuracy of participants' \textit{own}
decisions more than causal explanations did. They also found the effect of familiarity
to be conditional on whether the AI was correct. Their sample size is a useful reference
point when interpreting null moderation results obtained on smaller samples.

\subsection{Research Positioning}

Table~\ref{tab:positioning} situates this study against its three nearest antecedents.
The present work is distinguished by the conjunction of three properties: explanation
\textit{types} are compared against one another rather than against a no-explanation
baseline; the comparison occurs inside a multi-turn LLM-driven conversation rather than a
static display; and the domain is sensory-hedonic with a non-digitally-transmissible core
attribute. In addition, the narrative-based condition operationalises explanation
\textit{content} that is affective, which is distinct from the conversational
\textit{style} manipulations that dominate the anthropomorphism literature
\cite{libaoku2024}, \cite{cheng2022}.

\begin{table}[htbp]
\caption{Positioning Relative to Nearest Prior Work}
\begin{center}
\renewcommand{\arraystretch}{1.15}
\begin{tabular}{@{}p{0.9cm}p{1.6cm}p{1.5cm}p{1.5cm}p{1.5cm}@{}}
\toprule
& \textbf{Wang \& Benbasat \cite{wang2007}} & \textbf{Chen et al. \cite{chen2024}} & \textbf{Ziegfeld et al. \cite{ziegfeld2025}} & \textbf{This study} \\
\midrule
Setting & Static agent & Static display & Multi-turn CRS & Multi-turn CRS (LLM) \\
Manipu\-lation & Explanation type (how/why/ trade-off) & Presence and type & Preference elicitation & Explanation type (A/B/C) \\
Domain & Utilitarian & Both, compared & Utilitarian & Hedonic (perfume) \\
Trust & Three dimensions & Aggregate & Not focal & Aggregate \\
Outcome & Type-specific trust effects & Domain moderates type & No subjective effects & Type affects trust and PI \\
\bottomrule
\end{tabular}
\label{tab:positioning}
\end{center}
\end{table}

% ===============================================================
% SECTIONS III-VII TO FOLLOW
% ===============================================================
\section{System Design}
\label{sec:system}
\PH{Section III to be written next.}

\section{Method}
\label{sec:method}
\PH{Section IV to be written next.}

\section{Results}
\label{sec:results}
\PH{Section V to be written next.}

\section{Discussion}
\label{sec:discussion}
\PH{Section VI to be written next.}

\section{Conclusion and Future Work}
\label{sec:conclusion}
\PH{Section VII to be written next.}

% ===============================================================
\begin{thebibliography}{00}

\bibitem{jannach2021}
D. Jannach, A. Manzoor, W. Cai, and L. Chen, ``A survey on conversational
recommender systems,'' \textit{ACM Comput. Surv.}, vol. 54, no. 4, Art. 105,
May 2021.

\bibitem{rong2024}
Y. Rong \textit{et al.}, ``Towards human-centered explainable AI: A survey of user
studies for model explanations,'' \textit{IEEE Trans. Pattern Anal. Mach. Intell.},
\PH{vol., no., pp. --- source PDF is the accepted preprint; confirm published
volume, issue and page range}, 2024.

\bibitem{tintarev2012}
N. Tintarev and J. Masthoff, ``Evaluating the effectiveness of explanations for
recommender systems: Methodological issues and empirical studies on the impact of
personalization,'' \textit{User Model. User-Adapt. Interact.}, vol. 22,
\PH{pp. --- page range not visible in source PDF; confirm}, 2012.

\bibitem{chen2024}
C. Chen, A. D. Tian, and R. Jiang, ``When post hoc explanation knocks: Consumer
responses to explainable AI recommendations,'' \textit{J. Interactive Marketing},
vol. 59, no. 3, pp. 234--250, 2024, doi: 10.1177/10949968231200221.

\bibitem{brito2023}
R. de Brito Duarte, F. Correia, P. Arriaga, and A. Paiva, ``AI trust: Can
explainable AI enhance warranted trust?,'' \textit{Human Behav. Emerg. Technol.},
vol. 2023, Art. 4637678, 12 pages, 2023.

\bibitem{naiseh2021}
M. Naiseh, D. Al-Thani, N. Jiang, and R. Ali, ``Explainable recommendation: When
design meets trust calibration,'' \textit{World Wide Web}, vol. 24, pp. 1857--1884,
2021, doi: 10.1007/s11280-021-00916-0.

\bibitem{wang2007}
W. Wang and I. Benbasat, ``Recommendation agents for electronic commerce: Effects of
explanation facilities on trusting beliefs,'' \textit{J. Manage. Inf. Syst.},
vol. 23, no. 4, pp. 217--246, 2007, doi: 10.2753/MIS0742-1222230410.

\bibitem{stepin2021}
I. Stepin, J. M. Alonso, A. Catala, and M. Pereira-Fari\~{n}a, ``A survey of
contrastive and counterfactual explanation generation methods for explainable
artificial intelligence,'' \textit{IEEE Access}, vol. 9,
\PH{pp. --- page range not visible in source PDF; confirm}, 2021,
doi: 10.1109/ACCESS.2021.3051315.

\bibitem{varga2024}
M. Varga and P. Albuquerque, ``The impact of negative reviews on online search and
purchase decisions,'' \textit{J. Marketing Res.}, vol. 61, no. 5, pp. 803--820,
2024, doi: 10.1177/00222437231190874.

\bibitem{wangx2022}
X. Wang, A. R. Ashraf, N. Thongpapanl, and K.-Y. Wang, ``Perceived deception and
online repurchase intention: The moderating effect of product type and consumer
regulatory orientation,'' \textit{J. Consumer Behav.},
\PH{vol., no., pp. --- source PDF is the early-view version; confirm final volume,
issue and page range}, 2022, doi: 10.1002/cb.2109.

\bibitem{libaoku2024}
B. Li, R. Yao, and Y. Nan, ``How does anthropomorphism promote consumer responses to
social chatbots: Mind perception perspective,'' \textit{Internet Research},
\PH{vol., no., pp. --- ahead-of-print at time of download; confirm final details},
2024, doi: 10.1108/INTR-04-2024-0583.

\bibitem{decampos2024}
L. M. de Campos, J. M. Fern\'{a}ndez-Luna, and J. F. Huete, ``An explainable
content-based approach for recommender systems: A case study in journal
recommendation for paper submission,'' \textit{User Model. User-Adapt. Interact.},
vol. 34, pp. 1431--1465, 2024, doi: 10.1007/s11257-024-09400-6.

\bibitem{ziegler2023}
D. C. Hernandez-Bocanegra and J. Ziegler, ``Explaining recommendations through
conversations: Dialog model and the effects of interface type and degree of
interactivity,'' \textit{ACM Trans. Interact. Intell. Syst.}, vol. 13, no. 2,
Art. 6, Apr. 2023, doi: 10.1145/3579541.

\bibitem{siro2023}
C. Siro, M. Aliannejadi, and M. de Rijke, ``Understanding and predicting user
satisfaction with conversational recommender systems,'' \textit{ACM Trans. Inf.
Syst.}, vol. 42, no. 2, Art. 55, Nov. 2023, doi: 10.1145/3624989.

\bibitem{rana2024}
A. Rana, S. Sanner, M. R. Bouadjenek, R. Di Carlantonio, and G. Farmaner, ``User
experience and the role of personalization in critiquing-based conversational
recommendation,'' \textit{ACM Trans. Web}, vol. 18, no. 4, Art. 43, Oct. 2024,
doi: 10.1145/3597499.

\bibitem{ziegfeld2025}
L. Ziegfeld, D. Di Scala, and A. H. M. Cremers, ``The effect of preference
elicitation methods on the user experience in conversational recommender systems,''
\textit{Comput. Speech Lang.}, vol. 89, Art. 101696, 2025,
doi: 10.1016/j.csl.2024.101696.

\bibitem{mcknight2002}
D. H. McKnight, V. Choudhury, and C. Kacmar, ``Developing and validating trust
measures for e-commerce: An integrative typology,'' \textit{Inf. Syst. Res.},
\PH{vol. 13, no. 3, pp. 334--359 --- not among the supplied source PDFs; verify
volume, pages and DOI independently}, 2002.

\bibitem{pu2012}
P. Pu, L. Chen, and R. Hu, ``Evaluating recommender systems from the user's
perspective: Survey of the state of the art,'' \textit{User Model. User-Adapt.
Interact.}, vol. 22, pp. 317--355, 2012, doi: 10.1007/s11257-011-9115-7.

\bibitem{davis1989}
F. D. Davis, ``Perceived usefulness, perceived ease of use, and user acceptance of
information technology,'' \textit{MIS Quart.},
\PH{vol. 13, no. 3, pp. 319--340 --- not among the supplied source PDFs; verify
independently}, 1989.

\bibitem{yen2021}
C. Yen and M.-C. Chiang, ``Trust me, if you can: A study on the factors that
influence consumers' purchase intention triggered by chatbots based on brain image
evidence and self-reported assessments,'' \textit{Behav. Inf. Technol.}, vol. 40,
no. 11, pp. 1177--1194, 2021, doi: 10.1080/0144929X.2020.1743362.

\bibitem{cheng2022}
X. Cheng, X. Zhang, J. Cohen, and J. Mou, ``Human vs. AI: Understanding the impact
of anthropomorphism on consumer response to chatbots from the perspective of trust
and relationship norms,'' \textit{Inf. Process. Manage.}, vol. 59, Art. 102940,
2022, doi: 10.1016/j.ipm.2022.102940.

\bibitem{liyang2025}
Y. Li, Z. Gan, and B. Zheng, ``How do artificial intelligence chatbots affect
customer purchase? Uncovering the dual pathways of anthropomorphism on service
evaluation,'' \textit{Inf. Syst. Frontiers}, vol. 27, pp. 283--300, 2025,
doi: 10.1007/s10796-023-10438-x.

\bibitem{crolic2022}
C. Crolic, F. Thomaz, R. Hadi, and A. T. Stephen, ``Blame the bot: Anthropomorphism
and anger in customer-chatbot interactions,'' \textit{J. Marketing}, vol. 86, no. 1,
pp. 132--148, 2022, doi: 10.1177/00222429211045687.

\bibitem{bayer2023}
S. Bayer, H. Gimpel, and M. Markgraf, ``The role of domain expertise in trusting and
following explainable AI decision support systems,'' \textit{J. Decision Syst.},
vol. 32, no. 1, pp. 110--138, 2023, doi: 10.1080/12460125.2021.1958505.

\bibitem{celar2023}
L. Celar and R. M. J. Byrne, ``How people reason with counterfactual and causal
explanations for artificial intelligence decisions in familiar and unfamiliar
domains,'' \textit{Memory \& Cognition}, vol. 51, pp. 1481--1496, 2023,
doi: 10.3758/s13421-023-01407-5.

\end{thebibliography}

\end{document}
